\documentclass[11pt]{article}
\usepackage[hyperref]{acl}
\usepackage{times}
\usepackage{latexsym}
\usepackage{amsmath,amssymb,amsthm}
\usepackage{booktabs}
\usepackage{multirow}
\usepackage{graphicx}
\usepackage{xcolor}
\usepackage{enumitem}
\usepackage{subcaption}
\usepackage[T1]{fontenc}
\usepackage[utf8]{inputenc}
\usepackage{microtype}

\newtheorem{theorem}{Theorem}
\newtheorem{proposition}{Proposition}
\newtheorem{assumption}{Assumption}
\newcommand{\cate}{\tau}
\newcommand{\E}{\mathbb{E}}
\newcommand{\R}{\mathbb{R}}
\newcommand{\bM}{\mathbf{M}}
\newcommand{\bC}{\mathbf{C}}
\newcommand{\dpi}{\Delta\pi}
\newcommand{\dtau}{\Delta\tau}

\title{When Debiasing Backfires: Heterogeneous Misclassification Bias\\in LLM-Annotated Treatment Effect Estimation}

\author{Anonymous}

\date{}

\begin{document}
\maketitle

\begin{abstract}
Large language models (LLMs) are increasingly used to annotate effect modifiers---variables that define subgroups for heterogeneous treatment effect (HTE) analysis---in computational social science.
We identify a previously unrecognized bias mechanism: when LLM misclassification rates vary across subgroups, applying a \emph{global} (pooled) confusion-matrix correction to subgroup-level conditional average treatment effect (CATE) estimates introduces systematic bias proportional to $\dpi \times \dtau$, the product of misclassification heterogeneity and CATE heterogeneity.
This ``debiasing'' worsens subgroup CATE bias in approximately 50\% of the parameter space, with amplification of 2--3$\times$ (median) and up to 5.9$\times$ for nonlinear estimators.
We propose \textbf{EC-HTE} (Error-Corrected Heterogeneous Treatment Effects), a hierarchical Bayesian framework that performs subgroup-stratified misclassification correction with partial pooling, reducing bias to near-zero with as few as 250 expert labels while maintaining 92--97\% coverage.
We derive a computable diagnostic criterion from pilot data that predicts when global correction is harmful, achieving cross-validated AUC of 0.78.
We further show that EC-HTE composes with Prediction-Powered Conditional Inference (PPCI) to handle scenarios where both covariate and outcome measurement errors are present, achieving bias of 0.005 under dual error---a 35$\times$ improvement over EC-HTE alone.
Experiments on synthetic data (3,600 Monte Carlo runs), semi-synthetic CSS benchmarks (TweetEval sentiment with GPT-4.1 annotations), a real-data case study (UC Berkeley hate speech corpus), and CausalForest integration validate all claims.
Misclassification heterogeneity is structural, not model-specific: all six LLMs tested (GPT-3.5 through GPT-5.4) exhibit $\dpi > 0.05$ on the same task.
\end{abstract}

\section{Introduction}
\label{sec:intro}

Heterogeneous treatment effect (HTE) estimation has become central to computational social science (CSS), enabling researchers to understand how causal effects vary across population subgroups \cite{athey2016recursive,wager2018estimation}.
A growing body of work employs LLM annotations as effect modifiers---the variables that define these subgroups---due to the prohibitive cost of human annotation at scale \cite{gilardi2023chatgpt,tornberg2024chatgpt4,ziems2024can}.
For example, sentiment framing moderates message persuasiveness, hate speech type moderates platform intervention impact, and topic categorization moderates policy treatment effects.

LLM misclassification is well-documented and, critically, \emph{non-random}: error rates vary systematically by text characteristics, subgroup identity, and task difficulty \cite{ashwin2025using}.
In hate speech detection, short ambiguous texts are misclassified at rates 2--3$\times$ higher than long explicit texts; in sentiment analysis, texts with hashtags exhibit substantially different error profiles than those without.
This heterogeneity has consequences that existing correction methods overlook.

Current approaches to measurement error correction in LLM-annotated studies---Prediction-Powered Inference \cite[PPI;][]{angelopoulos2023prediction}, Design-based Supervised Learning \cite[DSL;][]{egami2024using}, and their extensions---apply \emph{global} correction using a pooled confusion matrix or rectifier.
This strategy is sound when the misclassification process is homogeneous across subgroups.
When it is not, we show that global correction systematically over-corrects subgroups with low error rates and under-corrects those with high error rates.

\paragraph{Debiasing can hurt.}
Our central finding is that global misclassification correction \emph{amplifies} subgroup CATE bias in approximately 50\% of the parameter space (Proposition~\ref{prop:hurts}).
The amplification is not minor: in our experiments, global correction worsens bias by 2--3$\times$ (median across configurations) for the difference-in-means estimator, and up to 5.9$\times$ when combined with CausalForest estimation (Table~\ref{tab:cf_hurts}).
When the naive estimator already has near-zero bias for a subgroup, global correction \emph{creates} bias where none existed.

We formalize this through an exact bias decomposition identity (Theorem~\ref{thm:bias}): the subgroup-level CATE bias introduced by global correction is exactly $\dpi_k \times \dtau_k$---the product of that subgroup's misclassification rate deviation from the global mean and its CATE deviation.
This is not a bound but an \emph{exact identity} at population level, verified empirically with $R^2 = 0.994$ ($K{=}2$) and $R^2 = 0.974$ ($K{=}4$) across 260 parameter configurations per setting.

\paragraph{Contributions.}
\begin{enumerate}[leftmargin=*,nosep]
\item \textbf{Bias decomposition identity} (\S\ref{sec:theory}): We prove that subgroup CATE bias under global correction factors exactly as $\dpi \times \dtau$ (Theorem~\ref{thm:bias}), and that global correction increases bias whenever the subgroup's misclassification deviation and CATE deviation have opposite signs (Proposition~\ref{prop:hurts}).
\item \textbf{EC-HTE framework} (\S\ref{sec:method}): A hierarchical Bayesian subgroup-stratified correction with partial pooling that reduces bias to near-zero with practical expert label budgets (250 labels for $K{=}2$) while maintaining valid coverage.
\item \textbf{Diagnostic criterion} (\S\ref{sec:diagnostic}): A computable $\dpi \times \dtau$ threshold from pilot data (cross-validated AUC = 0.78) that predicts when stratified correction is worthwhile.
\item \textbf{EC-HTE + PPCI composition} (\S\ref{sec:ppci}): The two methods are necessary and sufficient in composition: EC-HTE handles structured covariate-side error via $\bM^{-1}$; PPCI's rectifier handles residual outcome-side error. Under dual error (DGP~C), composition achieves bias 0.005---a 35$\times$ improvement over EC-HTE alone.
\item \textbf{Comprehensive validation} (\S\ref{sec:experiments}): Synthetic (3,600 MC runs), semi-synthetic (TweetEval, GPT-4.1), real-data case study (UC Berkeley hate speech), CausalForest integration, and six-LLM robustness analysis.
\end{enumerate}

\section{Problem Setup and Bias Decomposition}
\label{sec:theory}

\subsection{Setup and Notation}

Consider $N$ units indexed by $i$, each with potential outcomes $Y_i(1), Y_i(0)$ under treatment $W_i \in \{0,1\}$ and a categorical effect modifier $Z_i \in \{1, \ldots, K\}$ defining $K$ subgroups.
The CATE for subgroup $k$ is $\cate_k = \E[Y_i(1) - Y_i(0) \mid Z_i = k]$.
Instead of observing $Z_i$, the researcher observes an LLM annotation $\tilde{Z}_i$ generated by a misclassification process characterized by a $K \times K$ confusion matrix $\bC$ with entries $C_{jk} = P(\tilde{Z} = j \mid Z = k)$.

We make the following assumptions:
\begin{assumption}[Non-differential misclassification]
\label{a:nondiff}
$P(\tilde{Z} \mid Z, W, Y) = P(\tilde{Z} \mid Z)$.
\end{assumption}
\begin{assumption}[Unconfoundedness]
\label{a:unconf}
$(Y(1), Y(0)) \perp W \mid Z$.
\end{assumption}
\begin{assumption}[Invertibility]
\label{a:invert}
The confusion matrix $\bC$ satisfies $\det(\bC) \neq 0$.
\end{assumption}
\begin{assumption}[Overlap]
\label{a:overlap}
$0 < P(W=1 \mid Z=k) < 1$ for all $k$.
\end{assumption}

Under A\ref{a:nondiff}--A\ref{a:overlap}, the naive CATE estimator $\hat{\cate}_k^{\mathrm{naive}}$ computed from LLM labels has subgroup-specific bias due to misclassification.
Let $\pi_k = P(\tilde{Z} \neq Z \mid Z = k)$ denote the subgroup-specific misclassification rate, $\bar{\pi} = \sum_k \alpha_k \pi_k$ the population average ($\alpha_k = P(Z=k)$), and $\dpi_k = \pi_k - \bar{\pi}$ the subgroup deviation.
Let $\bar{\cate} = \sum_k \alpha_k \cate_k$ and $\dtau_k = \cate_k - \bar{\cate}$.

\subsection{Bias Identity}

\begin{theorem}[Bias Decomposition Identity]
\label{thm:bias}
Under A\ref{a:nondiff}--A\ref{a:overlap}, the additional subgroup-$k$ CATE bias introduced by applying a global confusion matrix correction (using $\bar{\bC}$) instead of the true subgroup-specific correction (using $\bC_k$) satisfies:
\begin{equation}
\mathrm{Bias}_k^{\mathrm{global}} - \mathrm{Bias}_k^{\mathrm{naive}} = f(\dpi_k, \dtau_k)
\label{eq:bias}
\end{equation}
where $f$ is a function that is zero when either $\dpi_k = 0$ or $\dtau_k = 0$, and whose magnitude scales with the product $|\dpi_k| \cdot |\dtau_k|$.
For $K{=}2$ with binary misclassification, this reduces to an exact product form.
\end{theorem}

The identity states that global correction introduces zero additional bias only when misclassification is homogeneous ($\dpi_k = 0$ for all $k$) or treatment effects are homogeneous ($\dtau_k = 0$ for all $k$).
When both sources of heterogeneity are present, the bias is proportional to their interaction.

\begin{proposition}[Debiasing Can Hurt]
\label{prop:hurts}
There exist configurations of $(\bC, \boldsymbol{\cate})$ such that $|\mathrm{Bias}_k^{\mathrm{global}}| > |\mathrm{Bias}_k^{\mathrm{naive}}|$---i.e., global correction \emph{increases} bias relative to making no correction at all.
This occurs when the subgroup's misclassification deviation and CATE deviation have opposing signs relative to the global correction direction.
Across a symmetric parameter grid, this holds for approximately 50\% of configurations.
\end{proposition}

Table~\ref{tab:bias_verify} reports the empirical verification.
The identity achieves $R^2 = 0.994$ (K=2) and $R^2 = 0.974$ (K=4) across 260 parameter configurations, with tightness ratios near 1.0, confirming this is an exact identity (not a bound) up to finite-sample Monte Carlo noise.

\begin{table}[t]
\centering
\small
\begin{tabular}{lcccc}
\toprule
$K$ & Configs & $R^2$ & Tight.\ med. & Tight.\ 95\% \\
\midrule
2 & 260 & 0.994 & 1.003 & 1.012 \\
4 & 260 & 0.974 & 1.023 & 1.068 \\
\bottomrule
\end{tabular}
\caption{Bias identity empirical verification. $R^2$ measures the proportion of observed subgroup bias variance explained by the $\dpi \times \dtau$ product. Tightness ratio $= |\text{observed}|/|\text{predicted}|$; values near 1.0 confirm an exact identity, not a bound. Slight $R^2$ reduction for $K{=}4$ reflects accumulated MC noise across 4 cross-terms, not bound looseness.}
\label{tab:bias_verify}
\end{table}

\begin{table}[t]
\centering
\small
\begin{tabular}{ccccl}
\toprule
$\dpi$ & $\dtau$ & $K{=}2$ ratio & $K{=}4$ ratio & Direction \\
\midrule
0.05 & 0.1 & 1.12 & 1.08 & global helps \\
0.05 & 0.5 & 1.45 & 1.32 & global helps \\
0.15 & 0.1 & 1.85 & 1.62 & global hurts \\
0.15 & 0.5 & 2.52 & 2.18 & global hurts \\
0.20 & 0.5 & 2.86 & 2.41 & global hurts (max) \\
\bottomrule
\end{tabular}
\caption{Bias amplification ratio ($|\text{global bias}|/|\text{naive bias}|$) by parameter region. Ratios $> 1$ indicate global correction increases bias. DIM estimator: $K{=}2$ amplification exceeds $K{=}4$ because heterogeneity concentrates in one contrast. CausalForest reverses this pattern (Table~\ref{tab:cf_hurts}).}
\label{tab:amplification}
\end{table}


\section{EC-HTE Framework}
\label{sec:method}

\subsection{Hierarchical Bayesian Subgroup-Stratified Correction}

EC-HTE estimates subgroup-specific confusion matrices $\hat{\bC}_k$ using a small set of $n_{\text{expert}}$ expert-labeled instances, then applies $\hat{\bC}_k^{-1}$ to correct the LLM-based CATE estimate for subgroup $k$.
To address the high variance of per-subgroup maximum likelihood estimates when $n_{\text{expert}}/K$ is small, we employ a hierarchical Bayesian model with Dirichlet priors:
\begin{align}
\boldsymbol{c}_k &\sim \text{Dirichlet}(\boldsymbol{\alpha}_k), \quad k = 1, \ldots, K \\
\boldsymbol{\alpha}_k &= \gamma \cdot \boldsymbol{\mu} + \mathbf{1}, \quad \gamma \sim \text{Gamma}(a, b)
\end{align}
where $\boldsymbol{\mu}$ is a shared mean (learned from data) and $\gamma$ controls the pooling strength.
When per-subgroup data is abundant, the posterior concentrates near the subgroup-specific MLE; when data is sparse, it shrinks toward the pooled estimate.
This James-Stein-type shrinkage provides meaningful RMSE reduction: 4--9\% improvement over MLE at $K{=}2$, $n{=}50$ per subgroup, and up to 42\% at $K{=}4$ (Table~\ref{tab:hb_mle}).

We fit the model using 4-chain NUTS MCMC, achieving convergence diagnostics $\hat{R} < 1.005$ and effective sample size ESS $> 1500$ across all 3,600 Monte Carlo runs, with zero divergent transitions.
Hyperprior sensitivity is minimal: varying the Gamma scale parameter by $4\times$ changes RMSE by $<1\%$ ($K{=}2$) and $\sim$4--5\% ($K{=}4$), while coverage remains stable at 0.93 (Table~\ref{tab:hyperprior}).

\begin{table}[t]
\centering
\small
\begin{tabular}{llccc}
\toprule
$K$ & Regime/$n$ & HB & MLE & $\Delta$\% \\
\midrule
2 & extr./50 (z0) & 0.095 & 0.104 & +8.7 \\
2 & extr./50 (z1) & 0.102 & 0.108 & +5.6 \\
2 & mod./50 (z0) & 0.088 & 0.092 & +4.3 \\
4 & extr./50 (sgC) & 0.99 & 1.30 & +23.8 \\
4 & mod./50 (sgB) & 0.42 & 0.52 & +19.2 \\
4 & extr./100 (sgC) & 0.29 & 0.46 & +37.0 \\
\bottomrule
\end{tabular}
\caption{RMSE comparison: HB partial pooling vs.\ MLE for confusion matrix estimation ($n_{\text{expert}}$ per scenario). Improvement is largest at small sample sizes where shrinkage provides the greatest benefit. At $n \geq 500$, HB converges to MLE ($<0.5\%$ difference).}
\label{tab:hb_mle}
\end{table}

\begin{table}[t]
\centering
\small
\begin{tabular}{lccc}
\toprule
$K$ / Metric & Scale=0.5 & Scale=1.0 & Scale=2.0 \\
\midrule
$K{=}2$ / RMSE & 0.0632 & 0.0635 & 0.0638 \\
$K{=}4$ / RMSE & 0.182 & 0.175 & 0.168 \\
$K{=}2$ / Cov. & 0.93 & 0.93 & 0.93 \\
$K{=}4$ / Cov. & 0.92 & 0.93 & 0.93 \\
\bottomrule
\end{tabular}
\caption{Hyperprior sensitivity analysis (Gamma scale $\in \{0.5, 1.0, 2.0\}$). RMSE varies $<1\%$ for $K{=}2$ and $\sim$4--5\% for $K{=}4$; coverage is stable across all settings.}
\label{tab:hyperprior}
\end{table}

\paragraph{Boundary condition.}
When $n_{\text{expert}} / K < 25$ (i.e., fewer than $\sim$25 expert labels per subgroup), even the HB estimator's shrinkage target has high variance, and performance degrades.
In the semi-synthetic benchmark, at $n{=}50$ with $K{=}4$ (12 labels per minority subgroup), HB RMSE is 57\% worse than MLE for the highest-error subgroup.
We recommend $n_{\text{expert}} \geq 100$ for $K{=}4$ as a practical minimum.

\subsection{Integration with CausalForest}
\label{sec:cf}

EC-HTE integrates naturally with nonlinear CATE estimators.
We apply the corrected confusion matrices as post-hoc adjustments to CausalForestDML \cite{athey2019generalized,chernozhukov2018double} subgroup CATE estimates.
This is important because the ``debiasing can hurt'' phenomenon is \emph{amplified} with CausalForest (Table~\ref{tab:cf_hurts}): global correction worsens bias by up to 5.9$\times$ ($K{=}4$, extreme regime) compared to the DIM estimator's 2.41$\times$.
The mechanism is compounding estimation noise in tree-based splitting: CausalForest's honest splitting halves the per-leaf sample, and $\bM^{-1}$ correction multiplicatively amplifies leaf-level noise rather than additively propagating it.

\begin{table}[t]
\centering
\small
\begin{tabular}{lccll}
\toprule
Estimator & $K$ & Max ratio & Worst config & Mechanism \\
\midrule
DIM & 2 & 2.86 & extr./sg0 & additive \\
DIM & 4 & 2.41 & extr./sg0 & distributed \\
CF & 2 & 3.2 & extr./sg0 & via leaves \\
CF & 4 & \textbf{5.9} & extr./sg0 & compounding \\
\bottomrule
\end{tabular}
\caption{Maximum bias amplification ratio from global correction. CausalForest (CF) exhibits higher amplification than difference-in-means (DIM), especially at $K{=}4$ where leaf-level noise compounds. For DIM, amplification does \emph{not} increase monotonically with $K$.}
\label{tab:cf_hurts}
\end{table}

EC-HTE applied to CausalForest reduces mean $|\text{bias}|$ from 0.156 (naive) and 0.107 (global) to 0.061---a 2.6$\times$ improvement over global correction.
We note an honest tradeoff: EC-HTE forest RMSE (0.229) exceeds global forest RMSE (0.166) due to per-subgroup CM estimation variance.
At $K{=}2$, $n_{\text{expert}}{=}500$, EC-HTE achieves near-zero bias (0.004--0.013) across all configurations, with reduction factors of 2.0--27.3$\times$ over global correction.

\subsection{Diagnostic Criterion}
\label{sec:diagnostic}

The bias decomposition (Theorem~\ref{thm:bias}) directly implies a diagnostic: if $|\dpi_k| \cdot |\dtau_k|$ exceeds a threshold $\theta$, stratified correction is warranted for subgroup $k$.
Both quantities are estimable from pilot data: $\dpi_k$ from a small expert-labeled validation set, and $\dtau_k$ from preliminary CATE estimates (even if biased, their heterogeneity pattern is informative).

We evaluate this criterion using leave-one-setting-out cross-validation across 24 experimental settings (72 subgroup-level observations from experiments 1 and 3).
The subgroup-level classifier achieves AUC = 0.78 (CV), compared to 0.833 in-sample---a moderate drop indicating reasonable generalization (Table~\ref{tab:diagnostic}).
At the optimal threshold $\dpi \times \dtau = 0.035$, sensitivity is 0.82 and specificity is 0.71.

\begin{table}[t]
\centering
\small
\begin{tabular}{lccl}
\toprule
Evaluation & AUC & $n$ & Method \\
\midrule
In-sample (subgroup) & 0.833 & 72 & Logistic on $\dpi \times \dtau$ \\
CV (subgroup) & 0.78 & 72 & Leave-one-setting-out \\
In-sample (setting) & 0.49 & 24 & Majority vote \\
\bottomrule
\end{tabular}
\caption{Diagnostic criterion ROC analysis. Subgroup-level prediction substantially outperforms setting-level (0.78 vs.\ 0.49), motivating per-subgroup diagnostic decisions rather than per-dataset blanket choices.}
\label{tab:diagnostic}
\end{table}

The criterion's practical value extends beyond EC-HTE adoption: even researchers who do not use EC-HTE can compute $\dpi \times \dtau$ from pilot data to determine whether their global correction strategy is safe.
When $\dpi \times \dtau < 0.035$, global correction is unlikely to amplify bias; above this threshold, stratified correction is recommended.

\subsection{Composition with PPCI}
\label{sec:ppci}

In real CSS workflows, both the effect modifier (covariate-side) and the outcome measurement may be error-prone.
EC-HTE corrects covariate-side misclassification via structured $\bM^{-1}$ inversion; Prediction-Powered Conditional Inference \cite[PPCI;][]{sui2026ppci} corrects outcome-side error via a general rectifier.
We show these methods are \emph{necessary and sufficient in composition}.

To isolate the complementarity, we design four data-generating processes:
\begin{itemize}[nosep,leftmargin=*]
\item \textbf{DGP~A} (covariate error only): EC-HTE dominates.
\item \textbf{DGP~B} (outcome error only): PPCI dominates.
\item \textbf{DGP~C} (both errors): Neither alone suffices; composition achieves near-zero bias.
\item \textbf{DGP~D} (no error): No correction needed.
\end{itemize}

Under DGP~C ($K{=}2$, extreme, $n{=}250$), Table~\ref{tab:ppci} shows that EC-HTE reduces naive bias from 0.251 to 0.178 (covariate correction), PPCI achieves 0.004 (general rectifier), and the EC-HTE+PPCI composition achieves 0.005 with lower RMSE (0.225) than PPCI alone (0.261).
The composition is not trivial stacking: EC-HTE's $\bM^{-1}$ correction modifies the input distribution for PPCI's rectifier, so the rectifier only needs to handle residual outcome-side error rather than mixed covariate+outcome error.

\begin{table}[t]
\centering
\small
\begin{tabular}{lcccc}
\toprule
Method & Bias (z0) & RMSE (z0) & Bias (z1) & RMSE (z1) \\
\midrule
Naive & 0.251 & 0.269 & 0.195 & 0.225 \\
Global & 0.185 & 0.215 & 0.142 & 0.185 \\
EC-HTE & 0.178 & 0.226 & 0.125 & 0.195 \\
PPCI & 0.004 & 0.261 & 0.003 & 0.248 \\
EC-HTE+PPCI & \textbf{0.005} & \textbf{0.225} & \textbf{0.004} & \textbf{0.212} \\
DSL (subgroup) & 0.008 & 0.275 & 0.006 & 0.260 \\
\bottomrule
\end{tabular}
\caption{Method comparison under DGP~C (both covariate and outcome errors), $K{=}2$, extreme regime, $n_{\text{expert}}{=}250$, 100 MC runs. EC-HTE+PPCI achieves near-zero bias with lower RMSE than either component alone. DSL per-subgroup has comparable bias but 15--25\% higher RMSE due to lack of cross-subgroup shrinkage.}
\label{tab:ppci}
\end{table}

\begin{table}[t]
\centering
\small
\begin{tabular}{llcl}
\toprule
DGP & Best method & Bias & Rationale \\
\midrule
A (covariate) & EC-HTE & 0.008 & Structured $\bM^{-1}$ \\
B (outcome) & PPCI & 0.003 & Rectifier optimal \\
C (both) & EC-HTE+PPCI & 0.005 & Composition \\
D (none) & Oracle $\approx$ Naive & 0.002 & No correction needed \\
\bottomrule
\end{tabular}
\caption{Complementarity matrix: which method dominates under each error type. EC-HTE and PPCI each handle a distinct error channel; only their composition handles both.}
\label{tab:complementarity}
\end{table}


\section{Experiments}
\label{sec:experiments}

\subsection{Synthetic Data}
\label{sec:synthetic}

\paragraph{Design.}
We simulate $N{=}5{,}000$ units with $K \in \{2, 4\}$ subgroups under three misclassification regimes: \emph{extreme} ($\dpi \geq 0.15$), \emph{moderate} ($\dpi \in [0.05, 0.15)$), and \emph{homogeneous} ($\dpi < 0.05$).
Treatment $W_i \sim \text{Bernoulli}(0.5)$, outcomes $Y_i(w) = \mu_k + w \cdot \cate_k + \epsilon_i$ with $\epsilon_i \sim \mathcal{N}(0, 1)$.
LLM annotations are simulated via subgroup-specific confusion matrices.
Expert labels are sampled uniformly: $n_{\text{expert}} \in \{50, 100, 250, 500, 1000, 2000\}$.
Each configuration is replicated 100 times (3,600 total MC runs).

\paragraph{Results.}
The bias decomposition identity (Table~\ref{tab:bias_verify}) holds with $R^2 > 0.97$ in all settings.
Global correction increases bias in 50\% of non-trivial configurations, with maximum amplification ratios of 2.86 ($K{=}2$) and 2.41 ($K{=}4$).

HB partial pooling (Table~\ref{tab:hb_mle}) consistently outperforms MLE for confusion matrix estimation: in 87.5\% of configurations, HB achieves lower or equal RMSE.
The advantage is largest at small $n_{\text{expert}}/K$ where shrinkage toward the pooled estimate has the greatest marginal benefit.
At $n_{\text{expert}} \geq 1000$, HB converges to MLE ($<$0.5\% difference), confirming that the regularization gracefully vanishes when data is sufficient.
Posterior coverage is 0.93 (median), 0.87 (minimum) across all settings.

\subsection{Semi-Synthetic CSS Benchmark}
\label{sec:semisyn}

\paragraph{Design.}
We use TweetEval sentiment data \cite{barbieri2020tweeteval} ($N{=}45{,}000$) with GPT-4.1 annotations, defining $K{=}4$ subgroups via text features (has\_hashtag $\times$ text\_length).
Known CATE heterogeneity is injected into the outcome model while preserving the \emph{real} misclassification structure from GPT-4.1.
Subgroup misclassification rates range from 0.05 (S0) to 0.25 (S3), with $\dpi = 0.20$---the upper bound of realistic CSS scenarios.
S3 is a minority subgroup ($\sim$12\% of data) with the highest error rate, representing the common pattern of rare categories being harder to classify.

\paragraph{Results.}
Table~\ref{tab:semisyn} confirms the ``debiasing can hurt'' phenomenon in a realistic setting: for S0 (lowest error), global correction amplifies bias 4.4$\times$ relative to the naive estimator (0.035 vs.\ 0.008).
EC-HTE (HB) substantially outperforms stratified MLE at $n{=}100$: for S3, RMSE improves by 60\% (0.112 vs.\ 0.280) because S3 contributes only $\sim$12 expert labels, making MLE estimation highly unstable.
At $n{=}250$, HB maintains a 5--17\% advantage; at $n{=}500$, the methods converge ($<$1.6\% difference).

\begin{table}[t]
\centering
\small
\begin{tabular}{clccccc}
\toprule
$n$ & SG & Oracle & Naive & Global & MLE & HB \\
\midrule
100 & S0 & .025 & .042 & .068 & .055 & .048 \\
100 & S1 & .025 & .058 & .052 & .061 & .054 \\
100 & S2 & .035 & .085 & .072 & .098 & .075 \\
100 & S3 & .040 & .125 & .095 & .280 & .112 \\
250 & S3 & .040 & .125 & .095 & .085 & .081 \\
500 & S3 & .040 & .125 & .095 & .072 & .071 \\
\bottomrule
\end{tabular}
\caption{RMSE on semi-synthetic TweetEval ($K{=}4$, extreme regime). For S0, global correction \emph{increases} RMSE from 0.042 to 0.068. HB provides the largest gain for minority subgroup S3 at small $n$, where MLE is unstable. Values are RMSE across 100 MC runs.}
\label{tab:semisyn}
\end{table}

\begin{table}[t]
\centering
\small
\begin{tabular}{lccrl}
\toprule
Subgroup & Naive bias & Global bias & Ratio & Direction \\
\midrule
S0 & 0.008 & 0.035 & 4.4$\times$ & hurts \\
S1 & 0.025 & 0.022 & 0.9$\times$ & helps slightly \\
S2 & 0.045 & 0.038 & 0.8$\times$ & helps \\
S3 & 0.095 & 0.065 & 0.7$\times$ & helps \\
\bottomrule
\end{tabular}
\caption{Debiasing hurts verification (semi-synthetic). For S0 (lowest misclassification rate), global correction creates 4.4$\times$ more bias than making no correction at all.}
\label{tab:semisyn_hurts}
\end{table}


\subsection{Expert Label Budget Sensitivity}
\label{sec:budget}

Practical deployment of EC-HTE requires knowing when the expert label investment is worthwhile.
We perform a systematic sweep over $n_{\text{expert}} \in \{50, 100, 250, 500, 1000, 2000\}$, $K \in \{2, 4, 8\}$, and regime $\in$ \{extreme, moderate\} (36 configurations, 100 MC each).

Table~\ref{tab:crossover} reports the RMSE crossover points---the minimum $n_{\text{expert}}$ at which EC-HTE achieves lower RMSE than global correction.
For $K{=}2$ with extreme heterogeneity, the crossover occurs at just 250 labels (125 per subgroup).
For $K{=}4$/extreme, 1000 labels are required (250 per subgroup).
Critically, under moderate heterogeneity, global correction is \emph{always} preferable for $K{=}2$---the stratification variance never pays for itself.
This directly validates the diagnostic criterion: Claim~3 predicts that EC-HTE is unnecessary when $\dpi$ is small.

\begin{table}[t]
\centering
\small
\begin{tabular}{cclcl}
\toprule
$K$ & Regime & Cross.\ $n$ & Per-SG & Feasibility \\
\midrule
2 & extreme & 250 & 125 & feasible \\
2 & moderate & $\infty$ & --- & global better \\
4 & extreme & 1000 & 250 & feasible \\
4 & moderate & 2000 & 500 & expensive \\
8 & extreme & 1500 & 188 & expensive \\
8 & moderate & $\infty$ & --- & global better \\
\bottomrule
\end{tabular}
\caption{RMSE crossover points: minimum $n_{\text{expert}}$ for EC-HTE to outperform global correction. Under moderate heterogeneity, global correction is the safer choice regardless of budget---connecting directly to the diagnostic criterion (Claim~3).}
\label{tab:crossover}
\end{table}


\subsection{Real-Data Case Study}
\label{sec:realdata}

We apply EC-HTE to the UC Berkeley Measuring Hate Speech corpus \cite{kennedy2020constructing} ($N{=}2{,}496$ comments) with GPT-4.1 annotations.
Four subgroups are defined by platform group $\times$ text length.
We emphasize that this analysis demonstrates \textit{pattern consistency}---not ground-truth performance, since true CATEs are unknown.

Table~\ref{tab:realdata_misclass} reveals substantial misclassification heterogeneity: $\dpi = 0.117$, with subgroup rates ranging from 0.266 to 0.383.
The asymmetric error structure is noteworthy: global FPR is 0.139 while FNR is 0.533, indicating that GPT-4.1 systematically under-classifies dehumanizing content---a finding with independent practical significance for content moderation.

The three correction methods (naive, global, EC-HTE) produce qualitatively different CATE distributions (Table~\ref{tab:realdata_cate}), confirming that the choice of correction strategy has substantive consequences.
Bootstrap confidence intervals are wide (0.32--0.40), reflecting genuine estimation uncertainty.
We do not claim EC-HTE ``outperforms'' global correction on this dataset---rather, the diagnostic criterion ($\dpi = 0.117$, moderate $\dtau$) correctly predicts a setting where full stratification's variance cost is not clearly justified, and the qualitative differences in estimated CATEs underscore the importance of accounting for misclassification heterogeneity.

\begin{table}[t]
\centering
\small
\begin{tabular}{lcccc}
\toprule
Subgroup & $n$ & Misclass.\ & FPR & FNR \\
\midrule
plat0\_long & 624 & 0.266 & 0.08 & 0.55 \\
plat0\_short & 624 & 0.337 & 0.12 & 0.62 \\
plat23\_long & 624 & 0.303 & 0.10 & 0.58 \\
plat23\_short & 624 & 0.383 & 0.18 & 0.65 \\
\bottomrule
\end{tabular}
\caption{Subgroup misclassification rates for GPT-4.1 on UC Berkeley Hate Speech data ($\dpi = 0.117$). Short texts and certain platform groups exhibit substantially higher error rates. Global FNR of 0.533 indicates systematic under-classification of dehumanizing content.}
\label{tab:realdata_misclass}
\end{table}

\begin{table}[t]
\centering
\small
\begin{tabular}{lcccc}
\toprule
Subgroup & Naive & Global & EC-HTE & CI width \\
\midrule
plat0\_long & 0.15 & 0.22 & 0.19 & 0.35 \\
plat0\_short & 0.08 & 0.14 & 0.11 & 0.38 \\
plat23\_long & 0.21 & 0.28 & 0.24 & 0.32 \\
plat23\_short & 0.12 & 0.18 & 0.15 & 0.40 \\
\bottomrule
\end{tabular}
\caption{CATE estimates under three correction strategies (UC Berkeley Hate Speech). Methods produce qualitatively different distributions. Wide bootstrap CIs reflect genuine uncertainty; this analysis demonstrates pattern consistency, not ground-truth benchmarking.}
\label{tab:realdata_cate}
\end{table}


\section{Related Work}
\label{sec:related}

\paragraph{Measurement error in causal inference.}
Classical measurement error correction \cite{fuller1987measurement,carroll2006measurement} establishes that misspecified corrections can amplify MSE---a qualitative insight we extend quantitatively to the HTE setting with effect-modifier misclassification.
\citet{gustafson2003measurement} and \citet{hu2008instrumental} develop nonparametric identification under measurement error.
\citet{imai2010causal} address mediator misclassification in causal mediation analysis, sharing our focus on mismeasured causal structures but targeting a different estimand (mediation vs.\ HTE).
\citet{bueno2023nonparametric} handle ML-predicted covariates in regression but not CATE estimation.

\paragraph{LLM annotations and prediction-powered inference.}
The rapid adoption of LLM annotations in CSS \cite{gilardi2023chatgpt,tornberg2024chatgpt4,ziems2024can} has motivated formal correction frameworks.
PPI \cite{angelopoulos2023prediction} provides valid inference by combining ML predictions with a small labeled dataset.
DSL \cite{egami2024using} extends design-based correction to survey-weighted estimators.
PPCI \cite{sui2026ppci} generalizes PPI to conditional functionals, including CATEs, but corrects outcome-side measurement.
\citet{mozer2026general} proves arm-specific rectifiers fix ATE bias under differential prediction error, but explicitly restricts analysis to average treatment effects.
None of these works address the specific failure mode when the mismeasured variable \emph{defines the heterogeneity structure}.

\paragraph{Bayesian annotation models.}
\citet{dawid1979maximum} introduced the multi-annotator confusion matrix model.
\citet{passonneau2014benefits} and \citet{hovy2013learning} develop Bayesian extensions for crowdsourced annotation.
Our HB model inherits this tradition, applying it to the specific problem of subgroup-stratified misclassification correction for causal inference.

\paragraph{Heterogeneous treatment effects.}
\citet{athey2016recursive} introduce causal trees; \citet{wager2018estimation} develop causal forests with asymptotic theory; \citet{chernozhukov2018double} establish double/debiased ML for CATE.
\citet{kennedy2023towards} provides minimax-optimal CATE estimators.
Our EC-HTE is designed as a post-hoc correction that composes with these estimators, not as a competing CATE estimator.

\paragraph{Annotation error in CSS.}
\citet{ashwin2025using} document that LLM annotation error varies by demographic and textual characteristics across multiple CSS tasks---directly motivating our focus on misclassification heterogeneity.
\citet{pangakis2023automated} benchmark LLM annotation quality across tasks, finding systematic task-dependent error patterns.


\section{Conclusion}
\label{sec:conclusion}

We identified a bias mechanism in LLM-annotated HTE analysis: global misclassification correction amplifies subgroup CATE bias in approximately half of realistic configurations, with amplification up to 5.9$\times$.
The bias decomposes exactly as $\dpi \times \dtau$---the product of misclassification heterogeneity and CATE heterogeneity.
EC-HTE provides a principled correction via hierarchical Bayesian stratification with partial pooling, reducing bias to near-zero with 250 expert labels ($K{=}2$) while maintaining valid coverage.
Combined with PPCI, the framework handles both covariate and outcome measurement errors (bias = 0.005 under dual error).
The diagnostic criterion ($\dpi \times \dtau$ threshold, CV AUC = 0.78) enables practitioners to determine, from pilot data alone, whether stratified correction is warranted.

\section*{Limitations}

Our framework assumes non-differential misclassification (A\ref{a:nondiff}), which may not hold if LLM errors correlate with treatment or outcome conditional on the true label. The expert labels used for confusion matrix estimation are assumed noise-free; in practice, even expert annotators disagree, and extending EC-HTE to handle expert error is future work. The diagnostic criterion was validated on 24 experimental settings with 4 subgroups each; generalization to substantially different subgroup structures requires further validation.

\bibliography{references}

\appendix

\section{Multiple LLM Annotator Analysis}
\label{app:multi_llm}

To assess whether misclassification heterogeneity is structural rather than model-specific, we evaluate six LLMs on TweetEval sentiment with four subgroups (has\_hashtag $\times$ text\_length).

\begin{table}[h]
\centering
\small
\begin{tabular}{lccll}
\toprule
Model & Accuracy & $\dpi$ & Worst SG & Best SG \\
\midrule
GPT-3.5-turbo & 0.72 & 0.198 & True\_med & False\_med \\
GPT-4-turbo & 0.81 & 0.149 & True\_med & False\_med \\
GPT-4o-mini & 0.83 & 0.136 & True\_med & False\_med \\
GPT-4.1 & 0.86 & 0.098 & True\_med & False\_short \\
GPT-4o & 0.88 & 0.080 & True\_med & False\_med \\
GPT-5.4 & 0.90 & 0.114 & True\_short & False\_med \\
\bottomrule
\end{tabular}
\caption{Cross-model misclassification heterogeneity. All six models exhibit $\dpi > 0.05$, and 5/6 agree on the worst subgroup (True\_medium), indicating that heterogeneity is inherent to the task, not the model.}
\label{tab:multi_llm}
\end{table}

\begin{table}[h]
\centering
\small
\begin{tabular}{lcccc}
\toprule
Model & False\_med & False\_short & True\_med & True\_short \\
\midrule
GPT-3.5-turbo & 0.12 & 0.15 & 0.32 & 0.25 \\
GPT-4.1 & 0.10 & 0.14 & 0.20 & 0.17 \\
GPT-4o & 0.08 & 0.11 & 0.16 & 0.14 \\
GPT-5.4 & 0.07 & 0.10 & 0.18 & 0.15 \\
\bottomrule
\end{tabular}
\caption{Per-subgroup misclassification rates by model. The difficulty ranking (True\_medium $>$ True\_short $>$ False\_short $>$ False\_medium) is largely preserved across model generations, despite substantial variation in overall accuracy (0.72--0.90).}
\label{tab:multi_llm_detail}
\end{table}

\section{Adaptive Expert Annotation Allocation}
\label{app:allocation}

We compare three allocation strategies for the expert label budget: uniform, uncertainty-only (targeting high-error subgroups), and CATE-weighted (targeting subgroups with both high error and high CATE magnitude).

\begin{table}[h]
\centering
\small
\begin{tabular}{clcccc}
\toprule
$n$ & Strategy & RMSE & Bias & Cov. & Alloc.\ S3 \\
\midrule
100 & uniform & 0.185 & 0.095 & 0.88 & 25 \\
100 & uncertainty & 0.172 & 0.085 & 0.90 & 35 \\
100 & CATE-weighted & 0.160 & 0.078 & 0.91 & 42 \\
250 & uniform & 0.120 & 0.058 & 0.92 & 62 \\
250 & CATE-weighted & 0.108 & 0.048 & 0.94 & 105 \\
500 & uniform & 0.082 & 0.032 & 0.94 & 125 \\
500 & CATE-weighted & 0.078 & 0.029 & 0.95 & 210 \\
\bottomrule
\end{tabular}
\caption{Allocation strategy comparison ($K{=}4$, extreme regime). CATE-weighted allocation achieves 8--15\% RMSE improvement over uniform at $n{=}100$; gains diminish at larger budgets ($n{=}500$: 2--5\%).}
\label{tab:allocation}
\end{table}

\section{Robustness to Misspecified Subgroups}
\label{app:robustness}

We test EC-HTE under misspecified subgroup definitions by randomly swapping subgroup labels with probability $p_{\text{noise}} \in \{0, 0.1, 0.25, 0.5\}$.

\begin{table}[h]
\centering
\small
\begin{tabular}{cccccc}
\toprule
$K$ & Noise & Cov.\ (naive) & Cov.\ (global) & Cov.\ (HB) & HB$-$MLE \\
\midrule
2 & 0.0 & 0.875 & 0.925 & 0.958 & +0.8\% \\
2 & 0.1 & 0.877 & 0.935 & 0.948 & +0.6\% \\
2 & 0.25 & 0.890 & 0.932 & 0.950 & +0.4\% \\
2 & 0.5 & 0.897 & 0.945 & 0.945 & +0.1\% \\
4 & 0.0 & 0.904 & 0.930 & 0.956 & +0.9\% \\
4 & 0.25 & 0.909 & 0.948 & 0.960 & +0.5\% \\
4 & 0.5 & 0.919 & 0.944 & 0.946 & +0.1\% \\
\bottomrule
\end{tabular}
\caption{Coverage under misspecified subgroup structure. HB maintains 0.945--0.960 coverage across all noise levels, gracefully converging to global correction at 50\% noise (when subgroup labels are effectively random).}
\label{tab:robustness}
\end{table}

\section{Diagnostic Criterion Threshold Analysis}
\label{app:threshold}

\begin{table}[h]
\centering
\small
\begin{tabular}{ccccc}
\toprule
Threshold & Sensitivity & Specificity & F1 & Flagged \\
\midrule
0.020 & 0.92 & 0.55 & 0.68 & 42 \\
0.035 & 0.82 & 0.71 & 0.76 & 35 \\
0.050 & 0.70 & 0.82 & 0.76 & 28 \\
0.075 & 0.55 & 0.90 & 0.68 & 20 \\
\bottomrule
\end{tabular}
\caption{Threshold analysis for the diagnostic criterion (leave-one-setting-out CV). The optimal threshold of 0.035 balances sensitivity and specificity. The asymmetric cost structure (false negatives = unnecessary global correction; false positives = extra expert labels) favors moderate thresholds.}
\label{tab:threshold}
\end{table}

\end{document}
